%==============================================================================
% 十一、AI 工具使用声明（按官方规定置于参考文献之后）
%==============================================================================
\phantomsection
\section*{AI 工具使用声明}
\addcontentsline{toc}{section}{AI 工具使用声明}

按照竞赛关于人工智能工具使用的规定，本文声明如下。

所用工具：Hermes Agent（Nous Research），底层模型 deepseek-v4-flash（深度求索 DeepSeek），使用日期 2026 年 10 月\cite{deepseek}。工具与模型分属两家机构，此处分列以免混淆。

使用环节：代码编写与调试、图表绘制脚本的整理、文稿的文字润色。

未使用 AI 的环节：选题判断、建模思路与模型选型的确定、评价指标的选取、数据核验结论的判断、结果的取舍与结论提炼，均由参赛队员独立完成。

AI 生成内容的人工核验与修正：参赛队员对全部数值与结论逐条复核，发现并修正了以下问题。其一，初版把附件4 的连续强度直接当有序类别处理，且对预测与真值各做一次分位映射，导致加权 kappa 虚高；现改为只对真值定切点、预测沿用同一映射。其二，初版漏掉了附件8 感官名称表中 24 个感官项被标注为分类变量这一事实，现已按连续强度统一处理并在正文说明依据。其三，问题三初版把缺陷物阈值一律取 $\log_{10}$，使处于线性量纲域的双乙酰软约束退化为常数、$f_2$ 被固定在 0.524；现已逐列判断量纲域并修正，$f_2$ 可以真正降到 0。其四，初版引用的重采样次数与代码不符，现按实际运行结果改为问题一 1000 次、问题二 200 次。上述修正过程与结果均记录在支撑材料中，可复现。

参赛队对本文引用 AI 工具生成内容的正确性、可靠性与学术规范性负责。
